% LSPV-5 NOTE::SPARSEEX1::2026-10-02T08:03:00-07:00
\documentclass[11pt]{article}\usepackage[margin=.88in]{geometry}\usepackage{amsmath,amssymb}
\begin{document}
\section*{SPARSEEX1: exact exclusion of the monomial-plus-constant Boolean-grid exception}
\textbf{STATUS: PROVED/EXACT for the canonical prime-field grid \(H=\{0,\ldots,h-1\}\) under \(p>h(h-1)/2\). P versus NP remains OPEN.}
\paragraph{Literature audit.}
General sparse-polynomial root bounds over finite fields do not exclude many-root fewnomials: multiplicative subgroup/coset constructions give genuine exceptions. Interval/subgroup character-sum results control broader distribution questions but their standard short-interval thresholds do not directly reach the vMS regime \(h\ll p^{1/4}\). Thus the proof below uses the exact consecutive grid rather than importing an inapplicable equidistribution theorem.
\paragraph{Theorem.}
Let \(p\) be odd prime, \(2\le h<p\), and
\[
f(x)=ax^e+c\in\mathbb F_p[x],\qquad a\ne0,\quad1\le e\le h-1.
\]
Assume \(f(j)\in\{0,1\}\) for every \(j=0,\ldots,h-1\). If \(p>h(h-1)/2\), then \(f\) cannot be nonconstant.
\paragraph{Proof.}
Since \(f(0)=c\), \(c\in\{0,1\}\). If \(c=0\), then for every \(j=1,\ldots,h-1\), Booleanity gives \(aj^e\in\{0,1\}\), and nonzero \(aj^e\) forces \(aj^e=1\). Hence all \(h-1\) nonzero grid points are roots of \(ax^e-1\). If \(c=1\), then \(1+aj^e\in\{0,1\}\), and \(aj^e\ne0\), so \(1+aj^e=0\); all nonzero grid points are roots of \(ax^e+1\). In either case a degree-\(e\) binomial has \(h-1\) distinct roots, so \(e\ge h-1\). Since \(e\le h-1\),
\[
e=h-1.
\]
Thus for some \(C\ne0\),
\[
x^{h-1}-C=\prod_{j=1}^{h-1}(x-j).
\]
Comparing coefficients of \(x^{h-2}\),
\[
0=-\sum_{j=1}^{h-1}j=-\frac{h(h-1)}2\pmod p.
\]
Hence \(p\mid h(h-1)/2\), impossible when \(p>h(h-1)/2\). QED.
\paragraph{Consequence for DIGCORR1.}
In the vMS parameter regime \(p\) is vastly larger than \(h^2\). Therefore every nonconstant canonical Boolean-grid leaf has at least two distinct active nonzero exponents. For every \(1\le\delta<m\), if also \(m(h-1)<p-1\), the DIGCORR1 degeneracy condition
\[
(p-1)\mid\delta(e-e')
\]
cannot hold for distinct active exponents \(e,e'<h\). Hence every two-frequency phase
\[
\alpha f(x)+\beta f(g^\delta x),\qquad\alpha\beta\ne0,
\]
retains a nonconstant term for every legal small lag.
\paragraph{Remaining gap.}
Nonconstant phase is sufficient to invoke appropriate complete polynomial exponential-sum bounds, subject to standard exceptional polynomial hypotheses, but translating those complex-valued estimates through exact bit-predicate Fourier expansions into \(\mathbb F_2\)-rank \(\omega(m)\) for the triangular supported operator remains OPEN.
\end{document}
% END_OF_VERIFIED_NOTE::LSPV-5::SPARSEEX1